%! TeX program = lualatex
\documentclass[../main.tex]{subfiles}

\begin{document}
\section{A collection of challenging problems}

\begin{exercise} \label{ex:midterm-practice-1}
  % (Midterm, Fall 2019)
  The logarithmic function \(f(x) = \log_{a}(x) + b\) whose graph passes through points \((1,2)\) and \((2,1)\) is 
  \begin{enumerate}[label=\Alph*.]
    \item \(\ln(e^{2}x)\)
    \item \(\log_{1/2}\left(\frac{x}{4}\right) + 1\)
    \item \(\log_{1/2}(x) + 2\)
    \item \(\log_{2}(x) - 1\)
    \item \(\log_{2}(x) + 2\)
  \end{enumerate}
  \begin{hint}
    This is similar to Problem 11, Quiz 1 on WebWork. A typical mistake is forgetting or misapplying logarithm rules.
  \end{hint}
\end{exercise}
\vfill

\begin{exercise} \label{ex:midterm-practice-2}
  % (Midterm, Fall 2019) 
  If \(f(x) = 2(\ln(x))^{3}\), then the inverse function \(f^{-1}(x)\) is 
  \begin{enumerate}[label=\Alph*.]
    \item \(e^{\sqrt[3]{x/2}}\)
    \item \(\sqrt[3]{e^{x/2}}\)
    \item \(2(e^{x})^{3}\)
    \item \(2e^{x/2}\)
    \item The inverse function does not exist.
  \end{enumerate}
  
  \begin{hint}
    The systematic approach to find the inverse function is to isolate for \(y\) in the equation \(x = f(y)\) which comes from applying \(f(x)\) to both sides of \(y = f^{-1}(x)\). \updatenote{``comes from by applying'' corrected to ``comes from applying''.}  
    
    A typical mistake is to misapply the logarithm rules and conclude that \((\ln(x))^{3} = 3\ln(x)\).
  \end{hint}
\end{exercise}
\vfill

\begin{exercise} \label{ex:midterm-practice-3}
  % (Midterm, Fall 2022)
  Evaluate \(\lim_{x \to 0^{-}} e^{\frac{1}{\sin(x)}}\).

  \begin{enumerate}[label=\Alph*.]
    \item \(-\infty\)
    \item \(0\)
    \item \(e\)
    \item \(\infty\)
    \item The limit does not exist.
  \end{enumerate}
  \begin{hint}
    The systematic approach to limits of composite functions is change of variable. \updatenote{``composition functions'' corrected to ``composite functions''.} 
    
    A typical mistake is to directly substitute \(0\) into the function without realizing that the limit approaches \(0\) from the left and conclude that \(1/0 = \infty\) and \(e^{\infty} = \infty\). \updatenote{``A typically mistake'' corrected to ``A typical mistake''.}
  \end{hint}
\end{exercise}
\vfill
\clearpage

\begin{exercise} \label{ex:midterm-exercise-4}
  % (Midterm, Fall 2022)
  Suppose \(f(x) = \frac{1}{x - 1}\). How many discontinuities does the function \(f(f(x))\) have?

  \begin{enumerate}[label=\Alph*.]
    \item \(f(x)\) is continuous everywhere.
    \item Exactly one.
    \item Exactly two.
    \item Exactly three.
    \item More than three.
  \end{enumerate}
  
  \begin{hint}
    The systematic approach is to simplify a complex and unfamiliar expression. Notice that
    \[
      f(f(x)) = \frac{1}{\frac{1}{x-1} - 1} = \frac{1}{\frac{1-(x-1)}{x-1}} = \frac{x-1}{2-x}.
    \]
  \end{hint}
\end{exercise}
\vfill

\begin{exercise} \label{ex:midterm-exercise-5}
  % (Midterm, Fall 2022)
  Solve for \(x\) in the equation \(\ln(x) + \ln(x - 1) = 0\).
  \begin{enumerate}[label=\Alph*.]
    \item \(\frac{1 + \sqrt{5}}{2}\)
    \item \(\frac{1 + \sqrt{5}}{2}, \frac{1 - \sqrt{5}}{2}\)
    \item \(\frac{-1 + \sqrt{5}}{2}\)
    \item \(\frac{-1 - \sqrt{5}}{2}\)
    \item \(\frac{5}{2}\)
  \end{enumerate}
\end{exercise}
\vfill

\begin{exercise} \label{ex:midterm-exercise-6}
  % (Midterm, Fall 2022)
  Find the limit \(\lim_{h \to 0} \frac{(1 + h)^{100} - 4(1 + h)^{20} + 3}{h}\).
  \begin{enumerate}[label=\Alph*.]
    \item \(0\)
    \item \(-3\)
    \item \(20\)
    \item \(23\)
    \item The limit does not exist.
  \end{enumerate}
  
  \begin{hint}
    The systematic approach is to use the mind map on page~\pageref{page:mindmap-limits-zero-over-zero}. \updatenote{Added the missing ``the'' before ``mind map''.}  Problem 7 in Quiz 3 shows you how to recognize the derivative from a limit.
  \end{hint}
\end{exercise}
\vfill
\clearpage

\begin{exercise} \label{ex:midterm-exercise-7}
  % (Midterm, Fall 2023) 
  Find all values of \(c\) for which the function 
  \[
    f(x) = 
    \begin{cases}
      \frac{|x - 3|}{x - c}, &\text{if } x \ne c \\
      1, &\text{if } x = c
    \end{cases}
  \]
  is continuous on \((-\infty, \infty)\).
  \begin{enumerate}[label=\Alph*.]
    \item There are no such \(c\).
    \item \(c = 3\).
    \item All real numbers except \(3\).
    \item \(c = 0\).
    \item All real numbers except \(0\).
  \end{enumerate}
  
  \begin{hint}
    This is a very complex problem. The mind map on page~35 essentially guides us through the whole problem. \updatenote{``throught'' corrected to ``through''.}
    
    \begin{itemize}
      \item Start at the \enquote{continuity at \(a\)} node. \updatenote{Closing quote moved before ``node'' so the quote matches the node name in the mind map.}

       Systematically, go down the list of three tasks in the definition of continuity on page~\pageref{def:continuity}.  The first two tasks are satisfied.  The third task gives the equation \updatenote{``an equation involving'' corrected to ``the equation''; the sentence ran straight into the displayed equation.}
        \[
          \lim_{x \to c} f(x) = f(c).
        \]
      
      \item Write down the concrete expressions for both sides to see what we are dealing with.
        \[
          \lim_{x \to c} \frac{|x-3|}{x-c} = 1.
        \]

      \item In all examples (TODO insert ref) of such a problem, we attempted to evaluate the limit. We do the same here starting at the \enquote{try direct substitution} node in the mind map. \updatenote{``starting the'' corrected to ``starting at the'', and ``try direction substitution'' to ``try direct substitution''.}
        
    There are two possibilities
    \begin{itemize}
      \item If \(c ≠ 3\), then we have \(\frac{\text{nonzero}}{0}\) which does not exist and \(f(x)\) is not continuous at \(c\). So \(c\) must be \(3\). \updatenote{``So \(c\) cannot be \(3\)'' corrected to ``So \(c\) must be \(3\)''; this case rules out every \(c \ne 3\), not \(c = 3\).}
      \item If \(c = 3\), then we have a concrete equation
      \[
        \lim_{x \to 3} \frac{|x-3|}{x-3} = 1.
      \]
      We can directly evaluate the limit on the left-hand side using one-sided limits. The slightly faster way to do so is to recognize that the limit on the left-hand side is similar to 
      \begin{itemize}
        \item Part (b) of Problem 4 in Quiz 2,
        \item Part (d) of Problem 4 in Quiz 2, and
        \item Example~\ref{eg:derivative-absolute-value} on page~\pageref{eg:derivative-absolute-value}, where the derivative of \(|x-1|\) at \(1\) is the limit \(\lim_{h \to 0} \frac{|h|}{h}\).
      \end{itemize}
      
      Being able to recall these results helps you to quickly conclude that the limit on the left-hand side does not exist.  So \(c\) cannot be \(3\) either. \updatenote{``it to recognize'' corrected to ``is to recognize'', ``help you'' to ``helps you'', and ``cannot be anything other than \(3\)'' to ``cannot be \(3\)''; this case rules out \(c = 3\).}
    \end{itemize}
    \end{itemize}

    Finally, we conclude that there is no choice of \(c\) for which the function \(f(x)\) is continuous at \(c\).  
  \end{hint}
\end{exercise}
\vfill

\begin{exercise} \label{ex:midterm-exercise-8}
  % (Midterm, Fall 2023)
  If \(\lim_{x \to 8^{-}} f(x) = 1\), \(\lim_{x \to 8^{+}} f(x) = 1\), and \(f(8) = -1\), then which of the following statements is true?
  \begin{enumerate}[label=\Alph*.]
    \item \(f\) is continuous at \(8\).
    \item \(\lim_{x \to 8} f(x)\) exists and \(f\) is discontinuous at \(8\).
    \item \(\lim_{x \to 8} f(x)\) does not exist and \(f\) is discontinuous at \(8\).
    \item \(f\) is differentiable at \(8\).
    \item \(f'(8) = 0\).
  \end{enumerate}
  
  \begin{hint}
    The systematic approach is to sketch the given information as we did in Example~\ref{eg:discontinuity-classify-abstract} on page~\pageref{eg:discontinuity-classify-abstract}.
  \end{hint}
\end{exercise}
\vfill

\begin{exercise} \label{ex:midterm-practice-11}
  % (Midterm, Fall 2023)
  Let \(f(x) = \begin{cases} x \sin\left(\frac{1}{x^{2}}\right) &\text{if } x < 0 \\ 2x &\text{if } x \ge 0\end{cases}\), then 
  \begin{enumerate}[label=\Alph*.]
    \item \(f\) is not defined at \(0\).
    \item \(f\) is defined at \(0\) but has no limit there.
    \item \(f\) is differentiable at \(0\) but not continuous there.
    \item \(f\) is differentiable at \(0\)
    \item \(f\) is continuous at \(0\) but not differentiable there.
  \end{enumerate}
\end{exercise}
\vfill
\clearpage

\printhints

% \begin{enumerate}[wide, label=\textbf{Sample solution for Exercise~\thechapter.\arabic*.}, itemsep={1cm}]
%   \item 
%     This is a ``find all ... that satisfy ...'' type of question as discussed on page~\pageref{sec:working-with-equations}. We need equations involving the unknowns \(a\) and \(b\). The points \(f(x)\) passes through give us two equations. \updatenote{"pass" corrected to "passes" to agree with the singular subject "$f(x)$".}
%     \[
%       \left\{
%         \begin{array}{rcl}
%           f(1) &=& 2\\
%           f(2) &=& 1
%         \end{array} 
%       \right.
%       \longrightarrow
%       \left\{
%         \begin{array}{rcl}
%           \log_{a}(1) + b &=& 2\\
%           \log_{a}(2) + b &=& 1
%         \end{array} 
%       \right.
%       \overset{\text{\tiny \(\log_{a}(1) = 0\)}}{\longrightarrow}
%       \left\{
%         \begin{array}{rcl}
%           b &=& 2\\
%           \log_{a}(2) + 2 &=& 1
%         \end{array} 
%       \right.
%       \overset{\text{\tiny see below}}{\longrightarrow}
%       \left\{
%         \begin{array}{rcl}
%           b &=& 2\\
%           a &=& 1/2
%         \end{array} 
%       \right.
%     \]
%     In the last step, we solve for \(a\) in \(\log_{a}(2) = -1\). The fast method is to raise both sides to base \(a\) to get \(2 = a^{-1}\) and then \(a = 1/2\).  The systematic method (adds just one step) is to change base to \(e\) to get \(\frac{\ln(2)}{\ln(a)} = -1\) which turns into \(-\ln(2) = \ln(a)\). Raise both sides to \(e\) to get \(a = 2^{-1}\).

%     Alternatively, if you prefer exponentials over logarithms, start by converting \(f(x) = \log_{a}(x) + b\) into exponentials by raising both sides to \(a\) to get \(a^{f(x)} = x a^{b}\). Notice it is quite tempting to make an algebra mistake here. Plug in the point \((1,2)\) to get \(a^{2} = 1 a^{b}\) which means \(b = 2\). Plug in the point \((2,1)\) and \(b = 2\) to get \(a^{1} = 2 a^{2}\). Manipulate the exponents to get \(a = 2^{-1}\).


%   \item 
%     This is a mix of the fundamentals of inverse function, the equivalent equation in particular,  and the inverse pair \(e^{x}\) and \(\ln(x)\).  Start by writing down the equivalent equation for \(y = f^{-1}(x)\), as discussed in Example~\ref{eg:find-inverse}, which is \(x = 2(\ln(y))^{3}\). Isolate for \(y\) as follows:
%     \[
%       \frac{x}{2} = \left( \ln(y) \right)^{3} \quad\longrightarrow\quad \sqrt[3]{x/2} = \ln(y) \quad\longrightarrow\quad e^{\sqrt[3]{x/2}} = y.
%     \]
%     It is quite easy to make the following mistake: \(\Big(\ln(x)\Big)^{3} = 3\ln(x)\). The correct logarithm law is \(\ln(x^{a}) = a\ln(x)\). Brackets matter! \updatenote{First displayed equation had "$(\ln(x))^3$", inconsistent with $x=2(\ln(y))^3$ above and $\ln(y)$ in the next step; corrected to $(\ln(y))^3$. (The "common mistake" remark's use of $\ln(x)$ is unrelated and intentionally generic.)}

%   \item 
%     This is about applying limit laws (Theorem~\ref{thm:limit-laws} on page~\pageref{thm:limit-laws}) with only abstract information. Example~\ref{eg:limit-law-abstract} on page~\pageref{eg:limit-law-abstract} is relevant.
%     \[
%     \lim_{x \to a} \frac{2f(x)}{\sqrt{g(x)} - 1} = \frac{2 {\displaystyle \lim_{x \to a} f(x)}}{\sqrt{{\displaystyle \lim_{x \to a} g(x)}} - 1} = \frac{2 \times 3}{\sqrt{100} - 1} = \frac{2}{3}.
%     \]


%   \item 
%     This exercise involves a mix of domain, function composition, finding discontinuities (a subtask of classifying continuities discussed in class on page~\pageref{technique:classify-discontinuities}). 

%     Carry out the composition to get a concrete expression \(\frac{1}{\frac{1}{x-1} - 1}\). We can see a discontinuity \(x = 1\) because the inner quotient cannot be \(0\) in the denominator.  The zeros of the denominator of the whole quotient give us the remaining discontinuities.  We solve for \(x\) in \(\frac{1}{x - 1} - 1 = 0\). Therefore, \(x = 1\) and \(x = 2\) are the only discontinuities.

%     We would have missed the discontinuity at \(x = 1\) had we immediately manipulated \(f(f(x))\) to \(\frac{x-1}{2-x}\).  Remember, some algebraic operations change the domain of the function. 

%   \item 
%     This exercise is about properties (rules) of exponential functions and logarithms.  It is entirely an exercise in remembering the rules correctly and remembering that \(e^{x}\) and \(\ln(x)\) are inverses of each other. They ``cancel out.'' See the top of page~\pageref{eq:log-formulas} for a refresher of logarithm rules.

%     Use the logarithm rule \(\ln(ab) = \ln(a) + \ln(b)\). \[\ln(x) + \ln(x - 1) = 0 \quad\longrightarrow\quad \ln(x(x-1)) = 0 \quad\longrightarrow\quad x^{2} - x = 1.\] In the last step, we raise both sides to base \(e\).  Finish the calculation using the quadratic formula or completing the square. Pick only the positive one because \(\ln(\cdots)\) does not accept negative numbers as inputs. 

%   \item 
%     This question is about strategic thinking. The form of the limit suggests that the given limit is the derivative of something. Connect the idea used in Example~\ref{eg:recognize-derivative-from-limit} and differentiation rules.

%     From a problem-solving point of view, this exercise is harder than the brevity of its solution suggests. Please do not be lured into thinking that ``oh I \emph{just} need to recognize the derivative.'' The situational awareness for \emph{just} recognizing something is the result of frequent and deliberate reflection of definitions, techniques and strategies.

%     This limit is the \hlwarn{definition} for \(f'(1)\) where \(f(x) = x^{100} - 4x^{20}\).  We can do this by a guess-and-check process as discussed in Example~\ref{eg:recognize-derivative-from-limit} on page~\pageref{eg:recognize-derivative-from-limit}. Calculate \(f'(x) = 100 x^{99} - 80 x^{19}\) and plug in \(1\) to get \(20\).

%   \item 
%     This is a ``find all ... that satisfy ...'' type of question as discussed on page~\pageref{sec:working-with-equations}.  The function is continuous everywhere except possibly at \(c\). The \hlwarn{definition} of continuity at \(c\) gives us an equation \(\lim_{x \to c} \frac{|x - 3|}{x - c} = f(c)\) involving the unknown \(c\). Following the discussion on page~\pageref{sec:working-with-equations}, we are going to draw a conclusion for the continuity at \(c\) based on this equation.  Recall the three tasks discussed immediately below the definition of continuity on page~\pageref{def:continuity}.

%     The right-hand side is \(1\).  The left-hand side requires two separate analysis.
%     \begin{enumerate}
%       \item Case \(c \ne 3\).  Then \(\lim_{x \to c} \frac{|x - 3|}{x - c}\) is of type \(\frac{\text{non-zero constant}}{0}\) which does not exist.  See Example~\ref{eg:limit-constant-over-0} for a reference.

%       \item Case \(c = 3\). Notice \(\lim_{x \to c} \frac{|x - 3|}{x - c}\) is of type \(\frac{0}{0}\). We learned on page~\pageref{eg:limit-cancellation} that such a type of limits requires more work. The piecewise nature of the function suggests the boxed statement on page~\pageref{sec:review-piecewise} and the technique for the second half of Exercise~\ref{ex:limit-piecewise-and-quotient} on page~\pageref{ex:limit-piecewise-and-quotient} are relevant. We perform the following sequence of calculations.  Start by applying the techniques discussed on page~\pageref{eq:abs} to write \(|x - 3|\) as a piecewise function.
%         \[
%           |x - 3| = 
%           \begin{cases}
%             x - 3 & \text{if } x \ge 3 \\
%             -(x - 3) & \text{if } x < 3.
%           \end{cases}
%         \]
%         As shown in the second half of Exercise~\ref{ex:limit-piecewise-and-quotient} on page~\pageref{ex:limit-piecewise-and-quotient}, divide by \(x - 3\) to get a concrete expression
%         \[
%           \frac{|x - 3|}{{\color{main}x-3}} 
%           =
%           \begin{cases}
%             \frac{x-3}{{\color{main}x-3}} &\text{if } x > 3 \\
%             \frac{-(x-3)}{{\color{main}x-3}} &\text{if } x < 3
%           \end{cases}
%           =
%           \begin{cases}
%             1 &\text{if } x > 3 \\
%             -1 &\text{if } x < 3
%           \end{cases}
%           .
%         \]

%         The limit approaches the boundary of the two branches. The two one-sided limits are quite obviously not equal. Therefore, the two-sided limit does not exist.
%     \end{enumerate}

%     The two-sided limit at \(c\) does not exist in either case.  There is no choice of \(c\) for which the equation \(\lim_{x \to c} f(x) = f(c)\) is true. Therefore, the function cannot be continuous at \(c\) regardless the choice of \(c\).

%   \item 
%     This is about Theorem~\ref{thm:limit-two-one-sides} on page~\pageref{thm:limit-two-one-sides} and the \hlwarn{definition} of continuity on page~\pageref{def:continuity}. The given left and right limits at \(8\) are equal. Therefore, \(\lim_{x \to 8} f(x)\) does exist. The \hlwarn{definition} for continuity at \(8\) requires \(\lim_{x \to 8} f(x) = f(8)\) to be true.  Therefore, the limit does exist but \(f\) is discontinuous at \(8\). 

%   \item 
%     This is about assumptions of limit laws (the red text in Theorem~\ref{thm:limit-laws} on page~\pageref{thm:limit-laws}) and the relation between one-sided limits and two-sided limits (Theorem~\ref{thm:limit-two-one-sides} on page~\pageref{thm:limit-two-one-sides}).

%     We did an iClicker poll in Week 3 about a sketchy calculation for the limit in Exercise~\ref{ex:limit-textbook-2.28}. The takeaway was that limit laws could not be applied because the assumptions of the limit lwas were not satisfied (one of the limits did not exist). In this exercise, the calculation \[\lim_{x \to 2} f(x) + 2f(-x) = \lim_{x \to 2} f(x) + 2 \lim_{x \to 2}f(-x) = \text{DNE} + \text{DNE} = \text{DNE}\] is invalid for the same reason. The only option left is to calculate left and right limits independently.

%     We can read the graph of \(f(-x)\) directly from the graph of \(f(x)\). If \(x\) approaches \(2\) from the left for \(f(x)\), then, on the graph of \(f(x)\), \(x\) approaches \(-2\) from the right for \(f(-x)\).  If \(x\) approaches \(2\) from the right for \(f(x)\), then, on the graph of \(f(x)\), \(x\) approachs \(-2\) from the left for \(f(-x)\).

%     Another (problem-solving) approach is to sketch the graph for \(f(-x)\) as shown below. Alternatively ... hold up the paper against the light, look at the graph of \(f(x)\) from the back side and mentally flip the signs on the horizontal axis. 

%     \begin{center}
%       \begin{tikzpicture}
%         \begin{axis}[basic curve, xmin=-5, xmax=5, ymin=-5, ymax=5, title={Graph of \(f(x)\)}]
%           \addplot[supp, very thick] coordinates {(-5, -4) (-2, -1)};
%           \addplot[very thick] coordinates {(-2,  2) ( 2, -4)};
%           \addplot[main, very thick] coordinates {( 2,  2) ( 5,  5)};
%           \draw[thick, fill=white] (axis cs:-2,-1) circle[radius=0.15];
%           \draw[thick, fill=black] (axis cs:-2, 2) circle[radius=0.15];
%           \draw[thick, fill=black] (axis cs: 2,-4) circle[radius=0.15];
%           \draw[thick, fill=white] (axis cs: 2, 2) circle[radius=0.15];
%         \end{axis}
%       \end{tikzpicture}
%       \begin{tikzpicture}
%         \begin{axis}[basic curve, xmin=-5, xmax=5, ymin=-5, ymax=5, title={Graph of \(f(-x)\)}]
%           \addplot[supp, very thick] coordinates {(+5, -4) (+2, -1)};
%           \addplot[very thick] coordinates {(+2,  2) (-2, -4)};
%           \addplot[main, very thick] coordinates {(-2,  2) (-5,  5)};
%           \draw[thick, fill=white] (axis cs:+2,-1) circle[radius=0.15];
%           \draw[thick, fill=black] (axis cs:+2, 2) circle[radius=0.15];
%           \draw[thick, fill=black] (axis cs:-2,-4) circle[radius=0.15];
%           \draw[thick, fill=white] (axis cs:-2, 2) circle[radius=0.15];
%         \end{axis}
%       \end{tikzpicture}
%     \end{center}

%     The left limit is \(\lim_{x \to 2^{-}} f(x) + 2f(-x) = -4 + 2 \times 2 = 0\).  The right limit is \(\lim_{x \to 2^{+}} f(x) + 2f(-x) = 2 + 2 \times (-1) = 0\).  The two one-sided limits agree and, consequently, the two-sided limit is the common value \(0\).

%   \item 
%     Example~\ref{eg:limit-law-abstract} is relevant.
%     The graph says \(\lim_{x \to 0} f(x) = -1\), and the limit of the denominator is not \(0\). We can directly apply limit laws.
%     \[
%       \frac{0^{2} + {\displaystyle \lim_{x \to 0} f(x)} + 3}{\sqrt[3]{3{\displaystyle \lim_{x \to 0}f(x) + 5}}} = \frac{-1 + 3}{\sqrt[3]{3(-1) + 5}} = \sqrt[3]{4}.
%     \]

%   \item 
%     Work your way down the list.
%     \begin{enumerate}[label=\Alph*.]
%       \item The function is defined at \(0\) in the second branch. 
%       \item The left limit is \(\lim_{x \to 0^{-}} x \sin(x^{-2}) = 0\) by the squeeze theorem because \(-x \ge x \sin(x^{-2}) \ge x\).  The right limit is \(0\). \updatenote{Upper bound corrected from $0$ to $-x$: since $-1\le\sin(x^{-2})\le 1$ and $x<0$, multiplying by $x$ flips the inequalities to $x \le x\sin(x^{-2}) \le -x$, not $x \le x \sin(x^{-2}) \le 0$ (which fails whenever $\sin(x^{-2})<0$). Both bounds $x,-x\to0$, so the squeeze conclusion is unaffected.}  The two-sided limit at \(0\) exists.
%     \end{enumerate}


%     The next three options are all about the derivative at \(0\). It cannot be C because continuity is a critical ingredient of differentiation as discussed on page~\pageref{thm:differentiability-implies-continuity}.  It comes down to either D or E. 

%     The number \(0\) sits at the boundary of the two branches. Therefore, we are forced to work with the \hlwarn{definition} of the derivative at \(0\) as two one-sided limits.
%     \[
%       \lim_{h \to 0^{-}} \frac{h\sin(1/h^{2})}{h} = \lim_{h \to 0^{-}} \sin(1/h^{2}) \overset{\text{sub } u = h^{2}}{=} \lim_{u \to 0^{+}} \sin(1/u).
%     \]
%     This is the oscillation discussed on page~\pageref{eg:limit-oscillation}. Therefore, the left limit does not exist.

%     No need to go any further for this particular exercise. The derivative at \(0\) does not exist by the \hlwarn{definition} of the derivative at \(0\).

%   \item 
%     The only information available for \(f(x)\) is an inequality.  That's a sign for the squeeze theorem as demonstrated in Exercise~\ref{ex:limit-squeeze-easy}.  The limit for the upper bound at \(0\) is calculated using the squeeze theorem again. However, this calculation is much simpler than Exercise~\ref{ex:limit-law-x-sin(x)} since the function \(|x \sin(1/x)|\) only outputs positive numbers.  The bounds are \(0 \le | x \sin(1/x)| \le |x|\). The limit of the upper bound function \(|x|\) as \(x \to 0\) is \(0\). Viola! The squeeze theorem says \(\lim_{x \to 0} f(x) = 0\).

%   \item 
%     This is about using a change of variable to turn an unfamiliar problem into a familiar one.
%     Substitute \(u = x^{1/6}\). As \(x \to 1\), \(u\) approaches the ``happy face'' \(\lim_{x \to 1} x^{1/6} = 1\). Moreover \(x^{1/3} = u^{2}\) and \(x^{1/2} = u^{3}\). After all substitutions, we arrive at
%     \[
%       \lim_{x \to 1} \frac{\sqrt[3]{x}-1}{\sqrt{x}-1} \overset{\text{sub \(u = x^{1/6}\)}}{=} \lim_{u \to 1} \frac{u^{2} - 1}{u^{3} - 1} = \lim_{u \to 1} \frac{\cancel{(u-1)}(u+1)}{\cancel{(u-1)}(u^{2}+u+1)} = \frac{2}{3}.
%     \]
% \end{enumerate}
\end{document}
